\scanpage{158}
$X$ schéma \uline{\unsure{Steinien}} sur $S$

\deleted{\illegible{1}} Soit $\underline{L}$ inversible \inserted{très ample} sur $X$, tel que $f_{*}(\underline{L})\inserted{=E}$ [loc. libre de rang $N$] \inserted{tel que} $i:X\to P(\deleted{\unsure{f_{*}\underline{L}}}E)$ soit défini.
% slanted
\margin{\unsure{inutile} ?}

A) Équivalence

(i) Donner une base de $E=f_{*}(\underline{L})$ i.e. un \uline{isom} $\underline{O}_{S}^{N}\xrightarrow{\sim}\deleted{\unsure{f_{*}(\underline{L})|E}}E$

(ii) Donner un isom $\varphi:P_{S}^{N}\xrightarrow{\sim}P(\deleted{\unsure{f_{*}\underline{L}}}E)$, $\oplus$ un isom $(\varphi^{-1}i)^{*}(\underline{O}(1))\simeq\underline{L}$

(iii) Donner un \deleted{\unsure{isomorphisme|isom}} morphisme
\[
j:X\longrightarrow P_{S}^{N}
\]
$\oplus$ un isomorphisme $\alpha:j^{*}(\underline{O}(1))\simeq\underline{L}$
% slanted
\margin{\unsure{Lors|Pour|Dans} la définition des \unsure{équival\textsuperscript{ces}|équival\textsuperscript{ts}|équivalence}, faire attention qu'on exige que le changement de base le conserve \deleted{\uuline{\unsure{Non|N.B.}}, puisque on fait cette hypothèse \unsure{dans} \unsure{l'un|d'un} \unsure{indépendant|indépen\textsuperscript{t}}} \inserted{\uline{morphisme spécial}}}
tels que l'hom correspondant
\[
g_{*}\underline{O}(1)\longrightarrow f_{*}(\underline{L})=E
\]
soit un isom. [condition qui ne dépend pas du choix de $\alpha$]
\margin{Compatibilité aux opérations de $\unsure{\operatorname{GL}(N)|\operatorname{Gl}|\Gamma(\underline{O}_{S}^{*})}$ et localisation.}

A') Équivalence

(i) Donner une \deleted{\unsure{base de}} \inserted{section} du faisceau des germes de bases de $E$ modulo $\underline{O}_{S}^{*}$

(ii) Donner un isom
\[
\varphi:\deleted{\unsure{X\leftarrow|X}}P_{S}^{N}\longrightarrow P(E)
\]

(iii) Donner un morphisme
\[
j:X\longrightarrow P
\]
tel que \quad $1^{\circ})\quad j^{*}(\underline{O}(1))\underset{S}{\simeq}\underline{L}$ \uuline{\unsure{relativement à} $S$}

$2^{\circ})$ \quad L'hom $g_{*}(\underline{O}(1))\xrightarrow{\sim}f_{*}(j^{*}(\underline{O}(1))$ un isom
\[
\underbrace{\hspace{.77\linewidth}}_{\text{morphisme spécial}}
\]
% archive sheet number: 157
